\documentclass[10pt]{article}

\usepackage[letterpaper,margin=0.72in]{geometry}
\usepackage{amsmath,amssymb,booktabs,array,xcolor,microtype}
\usepackage[T1]{fontenc}
\usepackage{lmodern}
\usepackage{enumitem}
\usepackage{fancyhdr}
\usepackage{titlesec}
\usepackage[hidelinks]{hyperref}

\definecolor{erdosblue}{HTML}{2E74B5}
\definecolor{lightblue}{HTML}{EAF5FB}
\definecolor{lightgray}{HTML}{F2F4F7}

\pagestyle{fancy}
\fancyhf{}
\lhead{\footnotesize Erd\H{o}s Institute Quantitative Finance Project --- Summer 2026}
\rhead{\footnotesize Yue Wu, Freda Zhang}
\cfoot{\footnotesize \thepage}
\setlength{\headheight}{13pt}

\titleformat{\section}{\large\bfseries\color{erdosblue}}{}{0pt}{}
\titlespacing*{\section}{0pt}{7pt}{3pt}
\setlength{\parindent}{0pt}
\setlength{\parskip}{3.5pt}
\setlist[itemize]{leftmargin=1.3em,itemsep=1pt,topsep=2pt}

\newcommand{\formula}[2]{%
  \colorbox{lightblue}{%
    \parbox{0.965\linewidth}{\textcolor{erdosblue}{\textbf{#1}}\quad #2}%
  }%
}

\begin{document}

{\LARGE\bfseries\color{erdosblue} Executive Summary}\par
\vspace{2pt}
{\large\bfseries Market-Signal-Enhanced Delta Hedging on Historical Equity Paths}\par
\vspace{5pt}

\colorbox{lightgray}{\parbox{0.965\linewidth}{%
\textbf{Research question.} Can freely available, point-in-time market signals
improve replication of standardized European-call payoffs relative to
fixed-volatility Black--Scholes hedging?}}

\section{Objective and scope}

We test whether information beyond historical SPY prices improves dynamic
hedging. The experiment uses observed SPY, VIX, volume, sector-return, and
short-rate data, while defining standardized European calls instead of claiming
access to historical option quotes. It therefore measures payoff replication on
realized market paths, not historical bid--ask execution, early exercise, or
contract-level volatility-surface fit. The option claim, premium convention,
rebalancing calendar, and transaction-cost model are identical across
strategies.

\section{Mathematical framework}

For a call with spot $S_t$, strike $K$, time to maturity $\tau$, risk-free rate
$r$, and hedge volatility $\sigma_t$, every strategy holds the Black--Scholes
delta

\formula{Black--Scholes delta}{%
$\displaystyle \Delta_t=N(d_{1,t}),\qquad
d_{1,t}=\frac{\log(S_t/K)+(r+\tfrac12\sigma_t^2)\tau}
{\sigma_t\sqrt{\tau}}.$}

The strategies differ only in the volatility supplied to this common engine:

\begin{itemize}
  \item \textbf{FV-BS:} trailing 21-day realized volatility observed at
  initiation and then held fixed;
  \item \textbf{Rolling-BS:} trailing 21-day realized volatility updated on
  every hedge date;
  \item \textbf{VIX-BS:} contemporaneous VIX, used as a freely available
  option-market volatility proxy; and
  \item \textbf{MSA-Delta:} $\sigma_t^{\mathrm{MSA}}=m(Z_t)\widehat\sigma_t$,
  where $\widehat\sigma_t$ is rolling realized volatility and $m(Z_t)$ depends
  on a market-risk state.
\end{itemize}

After standardizing all features with development-period means and standard
deviations, MSA-Delta uses

\formula{Risk score}{%
$\displaystyle R_t=0.30\,\mathrm{IVRank}_t+0.20\,\mathrm{VIX}_t
+0.20\frac{\mathrm{RV}_{21,t}}{\mathrm{RV}_{63,t}}
+0.10\,\mathrm{VolumeShock}_t+0.10\,\mathrm{Drawdown}_{63,t}
+0.05\,\mathrm{SectorDiv}_t+0.05\lvert\Delta\mathrm{Rate}_{21,t}\rvert.$}

Training-score quantiles determine the state and multiplier:

\formula{Risk-state rule}{%
$\displaystyle
Z_t=\begin{cases}
\mathrm{normal},&R_t<q_{70},\\
\mathrm{elevated},&q_{70}\le R_t<q_{90},\\
\mathrm{stress},&R_t\ge q_{90},
\end{cases}\qquad
m(Z_t)\in\{1.00,1.15,1.40\}.$}

\section{Empirical method}

The split is strictly chronological. Training (April 2016--August 2020)
estimates feature scaling and risk thresholds. Validation (August 2020--May
2022) selects among three predeclared MSA configurations using mean absolute
hedging error (MAE). The selected moderate rule is refit on training plus
validation data, and the held-out test (May 2022--December 2024) is evaluated
once. Every feature at date $t$ uses information available at or before $t$.

The test grid crosses maturities of 10, 21, and 42 trading days with strike/spot
ratios of 0.95, 1.00, and 1.05. Episodes begin every five trading days, yielding
1,149 claims per strategy. Hedges rebalance daily and pay 5 basis points of
proportional costs. Each strategy receives the same FV-BS premium within an
episode. Terminal replication error is

\formula{Terminal error}{%
$\displaystyle E_T=V_T-\max(S_T-K,0),$ where $V_T$ is the hedge portfolio after
stock liquidation and transaction costs.}

We report MAE, root mean squared error (RMSE), loss Expected Shortfall (ES),
turnover, and moving date-block bootstrap confidence intervals for paired MAE
improvements. Overlapping starts expand coverage; block resampling addresses
their time dependence.

\section{Held-out results}

\begin{center}
\renewcommand{\arraystretch}{1.13}
\begin{tabular}{lrrrr}
\toprule
\textbf{Strategy} & \textbf{MAE} & \textbf{RMSE} & \textbf{95\% loss ES}
& \textbf{Avg. cost} \\
\midrule
FV-BS      & \textbf{1.5523} & 2.5422 & 6.7463 & 0.4631 \\
Rolling-BS & 1.6074 & 2.4385 & 6.4496 & 0.4690 \\
VIX-BS     & 1.6956 & \textbf{2.3993} & \textbf{5.5171} & \textbf{0.4580} \\
MSA-Delta  & 1.6906 & 2.5368 & 6.5819 & 0.4662 \\
\bottomrule
\end{tabular}
\end{center}

FV-BS has the lowest MAE, so the test does not support a blanket claim that
additional signals improve average hedging accuracy. VIX-BS has the lowest
RMSE, 95\% loss ES, average cost, and turnover. The difference is economically
meaningful: forward-looking volatility reduces some large replication errors
even though it is less accurate under the average absolute-error criterion.

\section{Uncertainty and robustness}

Define paired improvement as
$I=\lvert E_{\mathrm{FV}}\rvert-\lvert E_{\mathrm{strategy}}\rvert$, so
$I>0$ favors the alternative. Moving-block bootstrap estimates are $-0.138$
for MSA-Delta (95\% CI $[-0.254,-0.039]$), $-0.057$ for Rolling-BS
($[-0.119,0.006]$), and $-0.148$ for VIX-BS ($[-0.255,-0.033]$). Thus neither
VIX-BS nor MSA-Delta improves MAE relative to FV-BS with statistical support.
The central ranking also survives transaction-cost assumptions of 0, 5, and 10
basis points: FV-BS retains the lowest MAE and VIX-BS the lowest RMSE. Results
are additionally reported for all nine maturity--moneyness combinations.

\section{Conclusion and limitations}

\colorbox{lightblue}{\parbox{0.965\linewidth}{%
\textbf{Main finding.} Forward-looking option-market information improves
tail-risk control, but additional composite-signal complexity does not
automatically produce a better hedge.}}

Fixed realized volatility is best for minimizing average absolute replication
error in this sample. VIX is the strongest simple benchmark for squared error
and adverse-tail performance. The validation-selected MSA rule does not add
robust value beyond the simpler alternatives. This negative result is useful:
signal complexity must be validated out of sample rather than justified only by
economic intuition.

VIX is an index-level, approximately 30-day volatility measure rather than
contract-level SPY implied volatility. The controlled European claims omit
historical option spreads, dividends, American early exercise, and
volatility-surface effects. Block-bootstrap inference mitigates, but does not
eliminate, dependence among overlapping episodes. The natural extension is to
repeat the same leakage-controlled experiment using historical contract-level
option quotes when licensed data become available.

\end{document}
